\section{第四条主线：从固定 workflow 走向自主探索}

\subsection{当前 Agent 仍然偏向模板化执行}
圆桌后段的一段讨论很有张力：嘉宾指出，今天很多 Agent 虽然已经会调用工具、会执行多步流程，但本质上仍然非常依赖固定 workflow 模板，最多只是出现了有限的局部涌现。真正理想中的 Agent，不是把既定流程执行得更快，而是在陌生问题面前能主动尝试多种路径、比较反馈、再决定下一步。

\textit{``摆脱固定的 workflow 走向自主探索''} 这句话几乎概括了 panel 对下一阶段能力的期待。

\begin{importantbox}{自主探索意味着什么}
如果用更工程化的语言改写，这里的 ``探索'' 至少包含三层能力：
\begin{itemize}
  \item 知道什么时候应该继续试，什么时候应该回退；
  \item 知道不同工具或不同路径的机会成本；
  \item 在不确定环境里仍能维护局部状态和目标函数。
\end{itemize}
\end{importantbox}

\subsection{2026 年的三个高概率方向}
panel 最后的展望可以被整理成三条可操作的判断。第一，\textbf{scaling of acting} 与 \textbf{scaling of environment} 会继续推进，工业界会持续堆 infra、数据和 action space。第二，研究者会尝试同时 scale reasoning、agent 与 environment，而不是孤立优化某一项。第三，Agent 会在更多行业进入 ``最后一公里''，于是个性化、安全和部署细节会成为新的研究热点。

\begin{table}[H]
\centering
\begin{tabular}{p{3.1cm}p{4.1cm}p{4.4cm}}
\toprule
\textbf{2026 方向} & \textbf{核心动作} & \textbf{难点} \\
\midrule
Scale acting & 扩大 tool space / action space & 训练效率与探索成本 \\
Scale environment & 构建更多真实或可复现环境 & 沙箱、复现、reward 设计 \\
Scale co-evolution & 同时优化 agent 与 environment 关系 & 方法统一性不足 \\
Last-mile deployment & 深入行业场景与企业系统 & 个性化、安全、权限、合规 \\
\bottomrule
\end{tabular}
\caption{Panel 末尾对 2026 年的判断，更像研究与产业共同的路线图}
\end{table}

\subsection{本章小结}
这场讨论的终点不是 ``Agent 已经成熟''，而是 ``Agent 刚刚从模板化工作流迈向自主探索的前夜''。未来一年最值得关注的，不只是模型更强，而是它们能否在更真实、更动态的环境中稳定完成开放问题。

